\chapter{Description de l'interface}
\minitoc

La console est une application web à page unique. La navigation se fait par la
barre latérale gauche, qui compte huit entrées. Les deux dernières — journal
d'audit et réglages — sont séparées des autres : les fonctions courantes en
haut, l'audit et la configuration en bas.

\section{Vue d'ensemble}

C'est l'écran d'accueil de l'exploitation. Il répond à la question « que s'est-il
passé depuis que je ne regardais pas~? ».

On y trouve, de haut en bas :

\begin{itemize}[leftmargin=*]
  \item les \textbf{décisions requises} — les gestes à effet durable qui
        attendent votre arbitrage. Ce bloc est traité en détail au
        chapitre~5~;
  \item les \textbf{fiches de qualification} — les menaces inédites que la
        plateforme a contenues sans savoir les nommer~;
  \item les \textbf{chiffres des vingt-quatre dernières heures} : interventions
        prises en charge, gestes appliqués, gestes annulés~;
  \item le \textbf{flux des actions exécutées}, du plus récent au plus ancien.
\end{itemize}

\begin{tcolorbox}[colback=green!4!white, colframe=darkgreen,
  title=\textbf{Bon à savoir --- le taux d'annulation}, fonttitle=\bfseries,
  boxrule=0.6pt]
C'est l'indicateur à surveiller en priorité. Il mesure la fréquence à laquelle
la plateforme doit défaire ce qu'elle vient de faire. Une valeur qui s'élève
durablement justifie de revoir les consignes de réponse, voire de rouvrir la
question de la posture d'autonomie.
\end{tcolorbox}

\section{Portefeuille}

La liste des interventions, \textbf{ordonnée par enjeu décroissant}. L'ordre
n'est pas chronologique : il combine six composantes explicables — gravité
constatée, importance de la machine touchée, nombre d'observations
concordantes, confiance moyenne des capteurs, fraîcheur, et deux parts issues
du catalogue métier : la priorité et la dangerosité.

Aucun poids n'est appris. Le décideur doit pouvoir justifier l'ordre du
portefeuille sans lire les paramètres d'un modèle.

Chaque ligne se déplie sur le détail de l'intervention : observations reçues,
décision prise, gestes appliqués, chronologie.

\section{Surveillance}

L'état de sécurité des plateformes déclarées sous surveillance. On y ajoute une
plateforme par la modale \textbf{Ajouter}, en indiquant son nom, son adresse et
sa criticité.

La carte situe géographiquement les plateformes surveillées autour du siège du
Centre~; le balayage radar signale celles dont l'état s'est dégradé.

\section{Réversibilité}

Le catalogue des trente-huit gestes que la plateforme sait appliquer. Pour
chacun : ce qu'il fait, son degré d'annulabilité, le geste inverse, ce qui
subsiste après annulation, et le nombre approximatif d'équipements touchés.

Cet écran est en lecture seule. Il sert à répondre à la question « qu'est-ce
que ce système est capable de faire~? » — et la réponse est un ensemble fini,
inspectable, publiable.

\begin{tcolorbox}[colback=green!4!white, colframe=darkgreen,
  title=\textbf{Bon à savoir}, fonttitle=\bfseries, boxrule=0.6pt]
Trois gestes du catalogue sont marqués \emph{irréversibles} : l'extinction
d'une machine, l'effacement d'un disque et la suppression d'un compte. La
plateforme ne les exécute \textbf{jamais} d'elle-même, quelles que soient les
circonstances.
\end{tcolorbox}

\section{Démonstration}

Le déclenchement des vingt-deux types du catalogue, seuls ou par famille, plus
un scénario de menace inédite. C'est l'écran des présentations et de la
recette.

\section{Rapports}

L'édition des rapports d'opérations. Rien n'est produit à l'ouverture de
l'écran : vous choisissez ce que le rapport doit couvrir, puis son format.
Le chapitre~7 y est consacré.

\section{Journal d'audit}

Le registre chaîné, consultable et filtrable par type d'événement ou par
intervention. Chaque ligne porte son numéro d'ordre, son horodatage, son type,
son auteur et son empreinte.

En tête de l'écran, la \textbf{vérification de la chaîne} : elle recalcule les
empreintes de bout en bout et indique si la chaîne est intacte.

\begin{tcolorbox}[colback=orange!5!white, colframe=orange!70!black,
  title=\textbf{Attention}, fonttitle=\bfseries, boxrule=0.6pt]
Si la vérification signale une chaîne rompue, il ne s'agit pas d'une anomalie
de fonctionnement mais d'un \textbf{incident de sécurité portant sur la
plateforme elle-même} : une entrée a été modifiée en dehors de l'application.
Préservez le fichier de base en l'état et ouvrez une enquête.
\end{tcolorbox}

\section{Réglages}

Les préférences de session — thème, affichage — et, pour l'administrateur, la
posture d'autonomie, la politique de réponse, le coupe-circuit et la gestion
des comptes. Le chapitre~6 détaille ces fonctions.

\section{L'assistant}

Une bulle présente sur tous les écrans ouvre l'assistant d'exploitation. Il
répond à des questions portant sur les données de la plateforme : « fais le
bilan du jour », « pourquoi cette intervention n'a-t-elle rien déclenché~? »,
« combien d'annulations cette semaine~? ».

\begin{tcolorbox}[colback=green!4!white, colframe=darkgreen,
  title=\textbf{Bon à savoir}, fonttitle=\bfseries, boxrule=0.6pt]
L'assistant ne dispose d'aucune information hors des données de la plateforme,
et ne fabrique jamais un chiffre. Une question hors de son périmètre reçoit un
refus explicite accompagné de ce qu'il sait faire — jamais une réponse inventée.
\end{tcolorbox}
